\chapter{双分量凝聚体的模型与对称性}
\label{ch:two-component-model}

\begin{learningobjectives}
完成本章后，读者应能：
\begin{enumerate}
  \item 区分均匀 Josephson/Rabi 耦合双组分模型与具有动量转移的 Raman 自旋轨道耦合模型；
  \item 由 Raman 单粒子哈密顿量求两条能带，并判断下带何时出现两个简并极小值；
  \item 从哈密顿量识别全局相位、相对相位、粒子数和平移对称性；
  \item 写出双分量接触相互作用的 TWA 场方程及正确 Weyl 修正；
  \item 判断参数淬火改变了哪些守恒量，以及哪些“条纹”结论不能由模型自动推出。
\end{enumerate}
\end{learningobjectives}

\section{最小双分量哈密顿量}

考虑准一维赝自旋分量 $L,R$，场算符满足
\begin{equation}
  \comm{\hat\Psi_\sigma(x)}{\hat\Psi_{\sigma'}^\dagger(x')}
  =\delta_{\sigma\sigma'}\delta(x-x').
\end{equation}
均匀线性耦合模型写成
\begin{align}
  \Hhat_0={}&\int\dd x\,
  \hat{\bm\Psi}^\dagger
  \left[h_0\identity_2+\frac{\delta}{2}\sigma_z-J\sigma_x\right]
  \hat{\bm\Psi},
  \qquad
  \hat{\bm\Psi}=\begin{pmatrix}\hat\Psi_L\\\hat\Psi_R\end{pmatrix},
  \label{eq:uniform-two-component-h0}\\
  \Hhat_{\rm int}={}&\int\dd x\left[
  \frac{g}{2}\left(
  \hat\Psi_L^{\dagger2}\hat\Psi_L^2+
  \hat\Psi_R^{\dagger2}\hat\Psi_R^2\right)
  +g_{12}\hat\Psi_L^\dagger\hat\Psi_R^\dagger
  \hat\Psi_R\hat\Psi_L\right],
  \label{eq:two-component-interaction}
\end{align}
其中 $h_0=-\hbar^2\partial_x^2/(2m)+V(x)$，$2J$ 是本书约定下单粒子对称与反对称态的能级劈裂，$\delta$ 是失谐。文献也常把 Rabi 频率记为 $\Omega=2J$；比较公式时必须先统一这一因子。

式\eqref{eq:uniform-two-component-h0}适用于内部态的零动量转移 Rabi 耦合，也适用于两条腿或两个局域模之间的 Josephson 耦合。它本身没有有限波矢，因而在均匀系统中不能单独选择一个非零条纹波数。

\section{Raman 耦合与自旋轨道项}

两束反向传播 Raman 光在实验室表象中产生带位置相位的耦合，例如
\begin{equation}
  -J\int\dd x\left(
  \ee^{2\ii k_Rx}\hat\Psi_L^\dagger\hat\Psi_R+{\rm h.c.}
  \right).
  \label{eq:raman-lab-coupling}
\end{equation}
通过自旋依赖规范变换去掉该相位后，单粒子哈密顿量可写成
\begin{equation}
  h_{\rm soc}=
  \frac{(p_x-\hbar k_R\sigma_z)^2}{2m}
  +\frac{\delta}{2}\sigma_z-J\sigma_x+V(x).
  \label{eq:raman-soc-hamiltonian}
\end{equation}
令 $E_R=\hbar^2k_R^2/(2m)$。在均匀空间中取 $p_x=\hbar k$，可把式\eqref{eq:raman-soc-hamiltonian}写成
\begin{equation}
  h_{\rm soc}(k)=\epsilon_0(k)\identity_2+d_z(k)\sigma_z-J\sigma_x,
  \quad
  \epsilon_0(k)=\frac{\hbar^2(k^2+k_R^2)}{2m},
  \quad
  d_z(k)=\frac{\delta}{2}-\frac{\hbar^2kk_R}{m}.
  \label{eq:soc-bloch-vector}
\end{equation}
因此两条单粒子能带为
\begin{equation}
  E_\pm(k)=\epsilon_0(k)\pm\sqrt{J^2+d_z(k)^2}.
  \label{eq:soc-single-particle-bands}
\end{equation}
这一步也固定了本书 $J$ 约定下的临界因子。零失谐时，$E_-(k)$ 在
\begin{equation}
  J<2E_R
  \label{eq:soc-double-minimum-condition}
\end{equation}
具有两个简并极小值
\begin{equation}
  k=\pm k_0,
  \qquad
  k_0=k_R\sqrt{1-\left(\frac{J}{2E_R}\right)^2};
  \label{eq:soc-minimum-wavevector}
\end{equation}
在 $J\geq2E_R$ 时两个极小值并合到 $k=0$。若采用实验文献常见的 $\Omega=2J$，同一条件写成 $\Omega<4E_R$。非零失谐一般使两个极小值不再简并，不能继续套用式\eqref{eq:soc-minimum-wavevector}。

设下带两个极小值的归一化自旋量为 $\chi_\pm$。相干叠加
\begin{equation}
  \bm\psi(x)=c_+\chi_+\ee^{\ii k_0x}
  +c_-\chi_-\ee^{-\ii k_0x}
  \label{eq:soc-two-minimum-superposition}
\end{equation}
在总密度中产生交叉项
\begin{equation}
  2\Re\!\left[c_+^*c_-\chi_+^\dagger\chi_-
  \ee^{-2\ii k_0x}\right].
  \label{eq:soc-stripe-overlap}
\end{equation}
故候选条纹波矢为 $Q=2k_0$，但振幅还取决于两个极小值是否同时宏观占据以及自旋量重叠 $\chi_+^\dagger\chi_-$。如果两自旋量正交，总密度交叉项消失，自旋密度仍可能调制。相互作用最终决定单极小值、双极小值相干叠加或其他相，而不是式\eqref{eq:soc-single-particle-bands}单独决定；这一模型结构和相图可参见实验与综述文献\cite{Lin2011,Li2012,Zhai2015}。

因此，后续章节将明确区分两条问题线：
\begin{enumerate}
  \item 式\eqref{eq:uniform-two-component-h0}的相对相位淬火、去相位与相分离动力学；
  \item 式\eqref{eq:raman-soc-hamiltonian}中由双极小值、相互作用和相干叠加共同形成的条纹候选态。
\end{enumerate}
均匀 $J$ 模型中的两分量相干性下降，不等同于 Raman 自旋轨道体系中条纹序的建立或消失。

\section{密度与自旋通道}

在平均场或正规排序层面定义
\begin{equation}
  n=n_L+n_R,
  \qquad
  s=n_L-n_R.
\end{equation}
对称耦合 $g_{LL}=g_{RR}=g$ 时，相互作用能密度可重写为
\begin{equation}
  \mathcal E_{\rm int}
  =\frac{g+g_{12}}{4}n^2
  +\frac{g-g_{12}}{4}s^2.
  \label{eq:density-spin-interaction}
\end{equation}
于是 $g+g_{12}$ 控制总密度通道的刚度，$g-g_{12}$ 控制自旋/相对密度通道。无 Rabi 耦合的均匀混合物在 $g_{12}>g$ 时对低波数自旋涨落不稳定，倾向于相分离；这并不自动意味着总密度形成周期条纹。Rabi 耦合、动能、有限尺寸与 SOC 带结构可以改变稳定边界，必须对具体模型求解。

写
\begin{equation}
  \psi_L=\sqrt{n_L}\ee^{\ii\theta_L},
  \qquad
  \psi_R=\sqrt{n_R}\ee^{\ii\theta_R},
  \qquad
  \phi=\theta_R-\theta_L,
\end{equation}
则均匀线性耦合能为
\begin{equation}
  \mathcal E_J=-2J\sqrt{n_Ln_R}\cos\phi.
  \label{eq:josephson-phase-energy}
\end{equation}
$J>0$ 把平衡相对相位钉扎在 $\phi=0$；相对密度 $s$ 是其共轭变量。大 $J$ 会增大相对相位涨落的恢复力，但“钉死”只是在低能近似中的说法，量子基态仍有有限宽度。

\section{对称性与守恒量账本}

\subsection{均匀线性耦合存在时}

全局变换 $\hat\Psi_{L,R}\to\ee^{\ii\chi}\hat\Psi_{L,R}$ 保持哈密顿量不变，总粒子数
$\hat N=\hat N_L+\hat N_R$ 守恒。若 $J\neq0$，独立的相对相位旋转不是对称性，$N_L-N_R$ 不守恒。均匀无陷阱时还具有连续平移对称和总动量守恒。

\subsection{淬火后关闭线性耦合}

若没有其他自旋翻转项，哈密顿量恢复
$U(1)_L\times U(1)_R$，$N_L$ 与 $N_R$ 分别守恒。相对相位不再被显式势能钉扎，但其空间梯度和相互作用仍产生动力学。一个瞬时淬火改变哈密顿量而不改变 $t=0$ 的量子态；对一般初始密度算符，新能量
\begin{equation}
  E_{\rm new}=\Tr(\rhohat_{\rm old}\Hhat_{\rm new})
  \label{eq:postquench-energy}
\end{equation}
随后守恒。纯态时才可把它写成相应的 bra--ket。它一般高于新哈密顿量基态能，差值是淬火注入能量，而不是与热浴交换的热量。

\subsection{Raman SOC 的对称性}

规范变换后的式\eqref{eq:raman-soc-hamiltonian}保持平移不变，但自旋与动量被锁定；实验室表象的普通平移要与 Raman 相位变换联合考虑。零失谐时还可能存在交换两个极小值的离散对称。条纹态涉及有限波矢密度相干，并可能破坏连续或由外场剩余的离散平移对称；具体对称群取决于 Raman 几何和边界。

\section{双分量 TWA 场方程}

在同一正交实空间格点上，双分量接触项的 Weyl 漂移并不完全相同。令
$\delta_{\mathcal C}=1/\Delta x$，则均匀耦合模型的闭合系统 TWA 为
\begin{align}
  \ii\hbar\partial_t\psi_L={}&\left[
  h_0+\frac{\delta}{2}
  +g(|\psi_L|^2-\delta_{\mathcal C})
  +g_{12}\left(|\psi_R|^2-\frac12\delta_{\mathcal C}\right)
  \right]\psi_L-J\psi_R,
  \label{eq:two-component-twa-l}\\
  \ii\hbar\partial_t\psi_R={}&\left[
  h_0-\frac{\delta}{2}
  +g(|\psi_R|^2-\delta_{\mathcal C})
  +g_{12}\left(|\psi_L|^2-\frac12\delta_{\mathcal C}\right)
  \right]\psi_R-J\psi_L.
  \label{eq:two-component-twa-r}
\end{align}
同分量四次项的漂移扣除一个投影基线；跨分量密度乘积的另一分量只扣除半个基线。这可由
\begin{equation}
  (\hat n_L\hat n_R)_W
  =\left(|\alpha_L|^2-\frac12\right)
   \left(|\alpha_R|^2-\frac12\right)
  \label{eq:cross-interaction-weyl}
\end{equation}
直接求导得到。双线性 $J$ 项是二次哈密顿量，不产生额外排序修正。

对 SOC 模型，记 $\bm\psi=(\psi_L,\psi_R)^T$，则两式可合写为
\begin{equation}
  \ii\hbar\partial_t\bm\psi
  =\left[h_{\rm soc}+\mathcal V_W(\bm\psi)\right]\bm\psi,
  \label{eq:soc-spinor-twa}
\end{equation}
其中 $\mathcal V_W$ 是由式\eqref{eq:two-component-twa-l}--\eqref{eq:two-component-twa-r}两个非线性对角元组成的矩阵。SOC 是二次单粒子项，不另产生 Weyl 修正；排序基线仍由实际保留的空间--自旋模式投影决定。

\subsection{可直接实现的 SOC 傅里叶传播器}

在均匀周期系统中，式\eqref{eq:soc-bloch-vector}给出每个 $k$ 上的精确单粒子步进
\begin{align}
  U_k(\Delta t)={}&\exp[-\ii\epsilon_0(k)\Delta t/\hbar]
  \left[
  \cos\!\left(\frac{\lambda_k\Delta t}{\hbar}\right)\identity_2
  -\ii\frac{\sin(\lambda_k\Delta t/\hbar)}{\lambda_k}
  \left(d_z(k)\sigma_z-J\sigma_x\right)
  \right],
  \label{eq:soc-fourier-propagator}\\
  \lambda_k={}&\sqrt{J^2+d_z(k)^2}.
\end{align}
$\lambda_k\to0$ 时应以解析极限或数值稳定的 $\operatorname{sinc}$ 实现第二项。一个 Strang 步可依次执行半个 $U_k$、完整的局域势与相互作用步、再半个 $U_k$。若局域步中还有不与密度矩阵对易的自旋翻转，需对相应 $2\times2$ 矩阵指数化，不能把两个分量各自乘相位。

规范变换也改变边界条件。实验室表象的场严格周期时，式\eqref{eq:raman-lab-coupling}在环上首先要求 $2k_RL\in2\pi\mathbb Z$；规范表象的两个自旋分量随后带相反的扭转因子 $\exp(\mp\ii k_RL)$。只有更强的条件 $k_RL\in2\pi\mathbb Z$ 成立时，两分量才同时适用普通周期 FFT；$k_RL$ 为奇数倍 $\pi$ 时则是反周期边界。一般情形应显式使用对应的扭曲动量网格，或直接在实验室表象传播式\eqref{eq:raman-lab-coupling}。初态噪声必须在同一个有限自旋--空间基中生成；若用 BdG 模抽样，正能模式、自旋量、粒子--空穴伙伴与传播器采用的规范和边界必须一致。

\section{无量纲控制参数}

均匀系统可用总密度 $n$ 和密度相互作用能
$E_n=n(g+g_{12})/2$ 建立尺度。常见无量纲量包括
\begin{equation}
  \gamma_s=\frac{g-g_{12}}{g+g_{12}},
  \qquad
  j=\frac{J}{E_n},
  \qquad
  \tilde\delta=\frac{\delta}{E_n},
  \qquad
  \kappa=\frac{\hbar^2k_R^2/(2m)}{E_n}.
\end{equation}
但真正的 TWA 风险还取决于每个截止模式的占据、系统长度、温度和淬火后演化时间。仅保持这些平均场比值相同，并不保证不同网格的量子涨落等价。

\section{模型建立检查表}

进入 BdG 或 TWA 前应逐项回答：
\begin{enumerate}
  \item $L,R$ 是内部态、空间双阱还是动量态？
  \item 线性耦合是否携带空间相位或动量转移？$J$ 与实验 Rabi 频率如何换算？
  \item $g,g_{12}$ 是哪一维度和哪种正规化约定下的参数？
  \item 淬火前后分别守恒哪些粒子数、动量和离散对称量？
  \item 条纹指总密度、自旋密度，还是释放后干涉图样？
  \item 候选有序相究竟破坏哪个 Hamilton 对称性？
\end{enumerate}

\begin{chaptercheck}
若只给出式\eqref{eq:uniform-two-component-h0}而没有 $k_R$、非局域相互作用、晶格或其他有限波矢机制，就不能把任何密度起伏自动解释成 Raman 条纹相。先从单粒子色散和线性稳定谱中证明存在目标波矢，再设计序参量。
\end{chaptercheck}

\section*{本章小结}

均匀 Josephson/Rabi 耦合与 Raman 自旋轨道耦合是不同模型：后者包含动量转移并可产生双极小值。$g_{12}>g$ 表示自旋通道趋向相分离，却不自动建立总密度条纹。淬火到 $J=0$ 恢复相对 $U(1)$ 对称并使两分量粒子数分别守恒。双分量 TWA 中，同分量和跨分量相互作用具有不同 Weyl 基线修正，必须在固定模式投影下实现。

\begin{exercises}
  \item 从 $n_L=(n+s)/2$、$n_R=(n-s)/2$ 推导式\eqref{eq:density-spin-interaction}。
  \item 对式\eqref{eq:uniform-two-component-h0}计算单粒子对称、反对称态能量并核对劈裂为 $2J$。
  \item 从式\eqref{eq:cross-interaction-weyl}推导式\eqref{eq:two-component-twa-l}中的跨分量修正。
  \item 列出 $J\neq0$ 与 $J=0$ 时所有连续对称性及对应守恒量。
  \item 说明实验室表象式\eqref{eq:raman-lab-coupling}怎样通过自旋依赖规范变换产生式\eqref{eq:raman-soc-hamiltonian}。
  \item 从式\eqref{eq:soc-single-particle-bands}求导，推导零失谐双极小值条件\eqref{eq:soc-double-minimum-condition}和位置\eqref{eq:soc-minimum-wavevector}。
  \item 证明式\eqref{eq:soc-fourier-propagator}是幺正矩阵，并实现 $\lambda_k\to0$ 时无除零误差的版本。
\end{exercises}
